\section{Research Question}
\label{sec:otel-weaver-exp-005-collector-to-intake-rq}

\subsection{Problem Statement}
\label{sec:otel-weaver-exp-005-collector-to-intake-rq-problem}

Process intelligence requires event logs that satisfy strict object-centric structural laws:
every event must be bound to a named process instance (\texttt{process.pi.instance\_id}),
every refusal must be sealed with a BLAKE3 witness hash, and every admitted span must carry
a provenance timestamp (\texttt{pi.intake.ingested\_at}) that anchors the event to the
collector's clock rather than the emitting service's clock. Production observability
infrastructure does not emit logs in this form. OpenTelemetry spans carry arbitrary
attributes chosen by application developers, with no guarantee that process-intelligence
namespace attributes are present, correctly typed, or consistently populated.

The problem EXP-005 addresses is: \textit{can a stock OTel Collector deployment, configured
with OTTL processors, serve as the sole gatekeeping and normalisation layer between arbitrary
live OTLP traffic and a process intelligence intake pipeline, without requiring any
modification to the emitting service?} If yes, the collector becomes the first law-enforcing
membrane and the upstream application stack remains observability-standard-compliant without
process-intelligence-specific instrumentation obligations.

\subsection{Old-Regime Assumption Challenged}
\label{sec:otel-weaver-exp-005-collector-to-intake-rq-old-regime}

The old regime assumes that process evidence is manufactured by instrumenting applications
with process-aware SDKs, custom middleware, or bespoke log formatters. Under this assumption,
the process intelligence layer is tightly coupled to the emitting service: changing the
process model requires changing the service instrumentation. This creates a coupling that
makes process intelligence expensive to retrofit and fragile under schema evolution.

EXP-005 challenges this assumption by positioning the OTel Collector as the law-enforcement
boundary. The collector's OTTL processor language is expressive enough to express attribute
presence checks (\texttt{filter/pi} dropping spans where
\texttt{attributes["process.pi.instance\_id"] == nil}), attribute transformations
(\texttt{transform/pi} lowercasing activity names and injecting ingestion timestamps), and
structured export (\texttt{file/pi\_json} writing OCEL~2.0 JSON). If the collector can
enforce these laws without service-side changes, the old-regime coupling assumption is
falsified for the class of services that already emit OTLP spans.

\subsection{Hypothesis}
\label{sec:otel-weaver-exp-005-collector-to-intake-rq-hypothesis}

A three-stage OTel Collector pipeline --- OTLP receiver, OTTL filter+transform processors,
file exporter --- is sufficient to gatekeep, normalise, and emit OCEL~2.0-structured JSON
from live OTLP traffic. Specifically: spans missing \texttt{process.pi.instance\_id} will be
silently dropped at the \texttt{filter/pi} stage; admitted spans will have activity names
lowercased and an ISO~8601 ingestion timestamp injected at the \texttt{transform/pi} stage;
and the \texttt{file/pi\_json} exporter will write structurally valid OCEL~2.0 JSON to
\texttt{output/pi\_intake\_events.json} on every span batch flush. The harness Span~A/B/C
taxonomy provides a falsifiable three-class test surface: Span~A must appear in the output,
Span~B must be absent, and Span~C must appear but be rejected by the downstream live-check
gate.

\subsection{Success Criteria}
\label{sec:otel-weaver-exp-005-collector-to-intake-rq-success}

\begin{enumerate}
  \item \texttt{output/pi\_intake\_events.json} is materialised by a live Docker collector
        run processing harness output, and contains exactly Span~A and Span~C records.
  \item Span~B is absent from \texttt{pi\_intake\_events.json}, confirming \texttt{filter/pi}
        dropped it at the namespace boundary.
  \item Every admitted span record carries \texttt{pi.intake.ingested\_at} (ISO~8601) and
        \texttt{pi.intake.collector\_version} fields injected by \texttt{transform/pi}.
  \item The \texttt{live-check-intake-receipt.json} BLAKE3 receipt is materialised at
        \texttt{receipts/} and matches the hash declared in
        \texttt{live-check-intake.manifest.yaml}.
  \item Span~C fails the downstream live-check gate (exp-003 \texttt{validator.rs}) due to
        absent \texttt{process.pi.witness.hash}, producing a sealed refusal record.
  \item All criteria verified and the checkpoint upgraded from PARTIAL to ALIVE.
\end{enumerate}
